\section{Two top words share paired residues}

This is a positive, local compatibility theorem.  It concerns one Haar
level, not a complete multilevel decoder.  In particular, it does not by
itself improve the worst-case head bound.

\paragraph{Statement.}
Fix distinct real prefix poles $p_1,\ldots,p_q$ and distinct exterior
query positions $D_1,\ldots,D_q$, with $D_i>\max_r p_r$.
For all sufficiently small $t>0$, every two sign words
$\sigma,\tau\in\{-1,1\}^q$ admit one common collection of paired
residues $(a_r,b_r)_{r=1}^q$ such that both words have their prescribed
strict signs under the top-level paired Haar reader.  No node bias is
needed.  The allowed smallness threshold for $t$ can depend only on the
fixed poles and queries, not on the two words.

\paragraph{The extra cubic feature.}
Writing $x=D_i-p_r$, the paired readout divided by its positive gap $t$
has the exact form
\[
 \frac{a_r}{x(x+t)}-\frac{b_r}{x(x-t)}-\frac{2}{x^2-t^2}
 =\frac{a_r-b_r-2}{x^2-t^2}
  -\frac{t(a_r+b_r)}{x(x^2-t^2)}.
\]
Thus it is $K_t\operatorname{diag}(w)v+
L_t\operatorname{diag}(w)u$ for the queried word $w$, where
\[
 (K_t)_{ir}=\frac1{(D_i-p_r)^2-t^2},\qquad
 (L_t)_{ir}=\frac1{(D_i-p_r)((D_i-p_r)^2-t^2)},
\]
and $v_r=a_r-b_r-2$, $u_r=-t(a_r+b_r)$ are independent real
parameters when $t\neq0$.  In particular, bounded $u_r$ can require
paired residue sums of order $1/t$.  The shared-leading-difference
obstruction in \texttt{JointFamilyBiasedCauchyPool.tex} therefore does
not apply to this two-feature system.

\paragraph{Interpolation on half the source points.}
At $t=0$, the columns associated with any source subset $Q$ are
\[
 \left[\frac1{(D_i-p_r)^2},\frac1{(D_i-p_r)^3}\right]_{r\in Q}.
\]
If $2|Q|\geq q$, this matrix has full row rank.  To verify strictness,
choose $q$ columns consisting of paired value/derivative columns and
possibly one unpaired value column.  Separate each paired pole to
$p_r,p_r+h_r$.  Borchardt's identity for the resulting square Cauchy
matrix $C$ says
\[
 \det(C\mathbin{\circ}C)=\det(C)\operatorname{perm}(C).
\]
Divide by the product of the within-pair gaps and let every $h_r$ tend
to zero.  The Cauchy determinant formula supplies a nonzero limit:
its only vanishing Vandermonde factors are exactly those within-pair
gaps.  The permanent has a strictly positive limit because all entries
are positive.  The derivative columns are twice the displayed cubic
columns.  Hence the selected minor is nonzero.  The same argument
without confluent columns shows that every principal square submatrix
of $K_0$ is nonsingular.  Continuity preserves these finitely many
nonzero minors for all sufficiently small positive $t$.

\paragraph{Simultaneous signs.}
Replacing the second word $\tau$ by $-\tau$ leaves its requested sign
inequalities unchanged, since its entire readout changes sign.
Therefore assume that the disagreement class
$Q=\{r:\sigma_r\neq\tau_r\}$ has at least $q/2$ elements.
Put $P=Q^{\mathrm c}$.

Choose coefficients supported on $P$, using only $K_t$, whose first-word
readout $f$ equals $\sigma_i$ for $i\in P$.  This is possible by the
principal-submatrix nonsingularity.  Those coefficients contribute the
same readout $f$ for the second word, since the source signs agree on
$P$.

By the full-row-rank property for $Q$, choose coefficients supported on
$Q$ whose first-word readout is
\[
 y_i=\begin{cases}
 0,&i\in P,\\
 \sigma_i(|f_i|+1),&i\in Q.
 \end{cases}
\]
Their second-word readout is $-y$.  The combined readouts are therefore
$f+y$ and $f-y$.  On $P$ both have the requested sign, with signed
margin one.  On $Q$,
\[
 \sigma_i(f_i+y_i)=\sigma_i f_i+|f_i|+1\geq1,\qquad
 \tau_i(f_i-y_i)=-\sigma_i f_i+|f_i|+1\geq1.
\]
This proves the claim, uniformly over the finite collection of pairs of
words.  All scalar denominators stay positive when
$0<t<\min_{i,r}(D_i-p_r)$.

\paragraph{Formal and written parts.}
The exact paired algebra is kernel-checked as
\texttt{jointFamily\_paired\_square\_cube\_identity}.
The support decomposition and simultaneous strict-sign conclusion are
kernel-checked as \texttt{jointFamily\_two\_word\_compatibility},
with the stated interpolation hypotheses.  The Cauchy rank argument
above is a written proof; it has not been imported as an axiom into
Lean.  Both declarations are in
\texttt{HeadComplexity/Results/JointFamilyPairedCompatibility.lean}.

\paragraph{The recursive gap.}
A full child decoder must also retain all lower-level inequalities.
The residues constructed here can diverge as $t$ tends to zero, so
their effect on those lower levels cannot be dismissed by continuity.
The checked split theorem in \texttt{JointFamilyHaarSplit.lean}
instead begins with compatible full child residues and makes bounded
perturbations.  Proving compatibility of two arbitrary full child
systems, or controlling the divergent higher-order corrections across
all levels, remains open.
